\newpage
\phantomsection
\label{prf:no-object-is-its-own-successor}

\begin{remark*}[Return]
\hyperref[thm:no-object-is-its-own-successor]{Return to Theorem}
\end{remark*}

\begin{theorem*}[No Object Is Its Own Successor]
Let $(P,S,1)$ be a Peano system. For every $n \in P$,
\[
S(n) \neq n.
\]
\end{theorem*}

\begin{proof}
\textbf{Professional Standard Proof.}~\\
Let
\[
D:=\{x\in P:S(x)\neq x\}.
\]
We prove that $D=P$ by induction.

First, $1\in D$. Indeed, if $S(1)=1$, then $1$ would be the successor of the
element $1\in P$, contradicting the Peano axiom that $1$ is not a successor.
Thus $S(1)\neq 1$.

Now assume $n\in D$. Then $S(n)\neq n$. If $S(S(n))=S(n)$, injectivity of
$S$ would imply $S(n)=n$, contradicting $n\in D$. Therefore
$S(S(n))\neq S(n)$, so $S(n)\in D$.

Thus $D$ contains $1$ and is closed under successor. By the induction
principle, $D=P$. Hence for every $n\in P$, $S(n)\neq n$.
\end{proof}

\begin{proof}
\textbf{Detailed Learning Proof.}~\\
Define
\[
D:=\{x\in P:S(x)\neq x\}.
\]
The goal is to prove that every element of $P$ lies in $D$.

For the base case, we show $1\in D$. This means showing
\[
S(1)\neq 1.
\]
If $S(1)=1$, then $1$ would be the successor of an element of $P$, namely
$1$ itself. This contradicts the Peano axiom that the distinguished element
$1$ is not a successor of any element of $P$. Therefore $S(1)\neq 1$, and so
$1\in D$.

For the inductive step, assume $n\in D$. Then
\[
S(n)\neq n.
\]
We must show that $S(n)\in D$, which means
\[
S(S(n))\neq S(n).
\]
Suppose instead that $S(S(n))=S(n)$. Since the successor map is injective,
this equality of successors would imply
\[
S(n)=n,
\]
contradicting the inductive hypothesis. Hence $S(S(n))\neq S(n)$, so
$S(n)\in D$.

Therefore $D$ is inductive. By the induction principle for Peano systems,
$D=P$. Thus every $n\in P$ satisfies $S(n)\neq n$.
\end{proof}

\begin{remark*}[Proof structure]
The proof uses induction on the subset of elements that are not equal to
their own successors. The base case uses the axiom that $1$ is not a
successor. The successor step uses the contrapositive form of injectivity:
from $S(n)\neq n$, one obtains $S(S(n))\neq S(n)$.
\end{remark*}

\begin{remark*}[Dependencies]~\\
\begin{itemize}
  \item \hyperref[ax:peano-one-not-successor]{Distinguished Element Is Not a Successor}
  \item \hyperref[ax:peano-successor-injective]{Injectivity of Successor}
  \item \hyperref[thm:induction-principle-for-peano-system]{Induction Principle for a Peano System}
\end{itemize}
\end{remark*}
